El sistema permite entrenar modelos al usuario utilizando el sistema definido en la sección \textbf{Parte 1}.

A partir de los parámetros definidos en configuración avanzada, se hace llamada al mismo método que se ha utilizado para los entrenamientos de los modelos del proyecto, y se comienza un nuevo ciclo de iteraciones de entrenamiento. Este nuevo modelo guardará sus \textit{checkpoints} en una carpeta a parte de la de los modelos entrenados previamente para no pisar el trabajo realizado.
Esto permite al usuario experimentar con el entrenamiento, ver cómo puede afectar la variación de hiperparámetros en los resultados y mostrar a alguien ajeno al tema qué es realmente un sistema de aprendizaje automático y , aunque no es un proceso instantáneo, y seguramente el usuario no llegue a esperar lo suficiente para obtener un modelo propio entrenado que ha aprendido lo suficiente como para generar resultados interesantes, se ha añadido esta funcionalidad porque se 
considera interesante que el usuario pueda ver en tiempo real cómo es realmente un entrenamiento de un modelo de inteligencia artificial.